%=====================================================================
\chapter{Literature Survey and Research Gaps}
\label{ch:litsurvey}
%=====================================================================

This chapter reviews the research directly relevant to the thesis. It is organised along the
same four axes as the research objectives, so that the gap identified at the end of each
section is the gap that the corresponding objective addresses. Emphasis is placed throughout on
what each body of work establishes and, more importantly, on what it leaves unresolved. Two
quantitative surveys conducted as part of this research --- one of the data types used in
recent multi-source studies and one of the fusion mechanisms they employ --- are reported in
Section~\ref{sec:fusionsurvey}, because the central gap of the thesis is visible in them
directly.

\section{Multi-Source Data Fusion in Financial Prediction}
\label{sec:fusionsurvey}

\subsection{That fusion helps is established}

The premise that heterogeneous sources outperform any single source is now well supported.
Snášel et al.~\cite{snasel2024} formalised a generalised multi-source fusion framework for stock
selection and demonstrated that complementary streams mitigate the individual weaknesses of each,
establishing the conceptual basis for combining numerical market data with unstructured text.
Long et al.~\cite{long2024} reached the same conclusion from a multi-view heterogeneous data
perspective, showing that different views of the same security carry non-overlapping
information. Nti, Adekoya and Weyori~\cite{nti2021} reported accuracy gains of 15--20 per cent
from a deep fusion of technical, fundamental and sentiment features, and --- a point that
recurs throughout this thesis --- emphasised that sources sampled at different frequencies must
be synchronised before they are combined, because misalignment can manufacture apparent
predictability where none exists.

The pattern recurs across a wide range of designs. Mu et al.~\cite{mu2023investor} combined
investor sentiment with optimised deep learning and reported consistent gains over price-only
baselines. Fu and Zhang~\cite{fu2024} examined news and social sentiment jointly and found that
the two behave differently: news sentiment supplies the more stable signal, social media the
more frequent directional one --- which is the empirical justification for treating them as two
streams rather than one. Shaban et al.~\cite{shaban2024} confirmed the pattern for trend
forecasting; Dong and Hao~\cite{dong2024} extended it to multi-dimensional, multi-level dynamic
fusion of macroeconomic and market data; Chen et al.~\cite{chen2024lasso} combined multi-feature
computation, penalised feature selection and a calendar-aware recurrent network to stabilise
performance in unsettled periods; Zhang et al.~\cite{zhang2018mil} aligned text with price
movements through multiple instance learning at appropriate time scales; Zhu~\cite{zhu2025informer}
processed multi-source social media data with a long-sequence transformer to capture long-range
patterns; Hong~\cite{hong2025mtl} predicted several correlated financial variables jointly
through a multi-task structure; and Alsheebah and Al-Fuhaidi~\cite{alsheebah2024} reported the
emerging-market case with endogenous and exogenous variables in a recurrent model.

Parallel work on single-source architectures establishes the ceiling that fusion must beat, and
it is a moving ceiling. Fischer and Krauss~\cite{fischer2018} established long short-term memory
as a strong baseline for financial market prediction on a large cross-section, and the
subsequent surveys of deep learning in financial time
series~\cite{sezer2020,ozbayoglu2020,jiang2021} document a decade of architectural refinement on
price-derived inputs. Song~\cite{song2024} enriched the feature space for trend-reversal
prediction and Javaid et al.~\cite{javaid2024} combined a peephole recurrent network with a
time-attention layer; both improve on conventional deep baselines while remaining confined to
price-derived inputs. A fusion architecture that beats an ordinary recurrent baseline has
therefore demonstrated little; the relevant comparison is against the strongest available
single-source model \emph{and} against a fixed-weight combination of the same sources, and
Chapter~\ref{ch:fusion} is constructed to supply both.

\subsection{A survey of data types and fusion mechanisms}

To locate the gap precisely rather than by impression, thirty-six recent multi-source studies
were surveyed and classified twice: by the primary data types they combine, and by the
mechanism through which they combine them. Table~\ref{tab:datatypes} reports the first
classification.

\begin{table}[!htb]
\centering
\caption[Distribution of data types in the surveyed multi-source literature]{Distribution of
primary data types across thirty-six recent multi-source financial prediction studies surveyed
as part of this research}
\label{tab:datatypes}
\tabfont
\begin{tabular}{@{}L{5.2cm}rrL{4.6cm}@{}}
\toprule
\textbf{Primary data category} & \textbf{Count} & \textbf{Share} &
\textbf{Characteristic limitation} \\
\midrule
Technical / OHLCV only & 9 & 25.0\% & Cannot access narrative information by construction \\
Textual or sentiment combined with technical & 16 & 44.4\% & Combination mechanism is almost
always static \\
Alternative data sources & 4 & 11.1\% & Feature extraction dominates; fusion is secondary \\
Reviews, multi-agent and other & 7 & 19.4\% & Do not report a directly comparable
architecture \\
\midrule
\textbf{Total} & \textbf{36} & \textbf{100\%} & \\
\bottomrule
\end{tabular}
\end{table}

The distribution shows a field that has resolved the question of \emph{what} to combine: a
plurality of recent studies already fuse text with market data, and a further tranche
incorporates alternative sources. Table~\ref{tab:fusiontax} reports the second classification,
and it shows that the question of \emph{how} to combine has not been resolved in the same way.

\begin{table}[!htb]
\centering
\caption[Taxonomy of fusion mechanisms in the surveyed literature]{Taxonomy of fusion
mechanisms in the surveyed multi-source studies, with the limitation each mechanism carries}
\label{tab:fusiontax}
\tabfont
\begin{tabular}{@{}L{2.6cm}L{4.1cm}L{4.9cm}L{2.2cm}@{}}
\toprule
\textbf{Mechanism} & \textbf{Description} & \textbf{Key limitation} &
\textbf{Representative work} \\
\midrule
Early fusion & Concatenation of feature vectors from all sources before modelling & Static
integration; assumes fixed inter-source relationships; prone to modality dominance &
\cite{moodi2024,mu2023investor,tai2022,li2023datafusion} \\
Late fusion & Averaging or stacking of predictions from source-specific models & Static,
pre-determined weighting; ignores time-varying source reliability &
\cite{farimani2024,vargas2022,alsheebah2024} \\
Attention-based hybrid & Neural attention weights features or modalities inside the network &
Weights are point estimates without uncertainty; can be unstable with noisy inputs &
\cite{liu2019attention,haryono2023,choi2023,zong2025} \\
Dynamic, uncertainty-aware & Weights derived from a time-varying estimate of each source's
predictive reliability & \textit{Not found in the surveyed literature} & --- \\
\bottomrule
\end{tabular}
\end{table}

The final row of Table~\ref{tab:fusiontax} is the finding. Across the surveyed studies, the
most sophisticated mechanisms in use learn a weighting, but they learn it once. Attention and
gating produce point estimates over an entire training set; normalising-flow
formulations~\cite{tai2022} model the joint distribution of sentiment and financial signals
without an explicit per-source confidence; gated cross-attention~\cite{zong2025} stabilises
multimodal fusion without measuring which modality has recently degraded. Farimani et
al.~\cite{farimani2024} describe their multimodal weighting as adaptive and move genuinely
toward time variation, but the adaptation is learned rather than measured, and no explicit,
probabilistic statement of source confidence enters the fusion. Bacco et al.~\cite{bacco2024}
study prediction explicitly under high uncertainty yet still fuse sentiment and price with a
static mechanism, which is precisely the configuration that should fail when uncertainty rises.

\subsection{Emerging markets and the Indian setting}

A further gap runs through the entire body of work reviewed above and is easiest to see when the
studies are sorted by market. The architectures are overwhelmingly designed, tuned and validated
on developed markets, and the calibrations they carry do not transfer cleanly.

Three instances make the point concrete. Yu and Liu's causal graph formulation~\cite{yu2025causalseqgnn}
is built on Fama--French factors~\cite{famafrench1993}, which are estimated on United States
equities and encode a size--value--profitability structure that does not describe an index in
which public-sector banks carry substantial weight, in which monsoon outcomes propagate through
agricultural commodities into the stocks exposed to them, and in which regulatory decisions by a
single domestic authority reset relationships across the board. Memory-augmented causal
networks~\cite{luo2025magicnet} store crisis patterns learned predominantly from United States
and Chinese data; the Indian market has had crises of its own --- the 2016 demonetisation shock
and the 2024 budget-related volatility among them --- whose dynamics do not resemble those
patterns. Multi-granularity graph constructions~\cite{zhu2025multigran} assume homogeneity across
instruments, which fails where highly liquid constituents and thinly traded ones behave
differently and where retail-driven noise is concentrated in the latter.

The specifically Indian evidence that does exist is narrow but pointed. Khan and
Chahal~\cite{khan2025asymmetric} establish, across eleven sectoral indices, that the influence of
social media sentiment on Indian sectoral returns is asymmetric and strongest at market extremes
--- a result that a fixed-weight fusion scheme cannot represent and that
Chapter~\ref{ch:fusion} recovers independently from its attention weights. Alsheebah and
Al-Fuhaidi~\cite{alsheebah2024} demonstrate the value of exogenous variables in an emerging-market
recurrent model. What is missing is a framework that jointly incorporates financial indicators,
news sentiment, social media signals \emph{and} Indian macroeconomic variables --- policy rate,
inflation, industrial production, currency --- under one architecture and one strict temporal
protocol, with the contribution of each stream measured rather than assumed. That absence is gap
G2 of Section~3.6 and the subject of Chapter~\ref{ch:fusion}.

\subsection{Fusion output rarely reaches a decision}

Where the output of a fusion architecture is carried through to a portfolio at all, the
conversion is typically heuristic. Sami et al.~\cite{sami2025} construct holdings by statistical
averaging of forecasts, without a covariance structure, position limits or an explicit risk
objective. The gap between a good forecast and a good portfolio is treated as an implementation
detail rather than as a modelling problem, and the consequence is that the reported economic
performance of such systems is not attributable to any identifiable component.

\section{Sentiment Analysis and Language Models in Finance}

\subsection{Three generations, and what each established}

Before the technical evolution, the substantive question: does text carry market-relevant
information at all? The affirmative evidence predates modern language models and is strong.
Tetlock~\cite{tetlock2007} showed that the pessimism of a widely read financial column forecasts
downward pressure on market prices followed by reversion, and later established with
co-authors~\cite{tetlock2008} that the linguistic content of news stories predicts firm earnings
and returns beyond the quantitative disclosure they accompany. Antweiler and
Frank~\cite{antweiler2004} found that the volume of internet message-board activity, not merely
its polarity, predicts volatility. Bollen, Mao and Zeng~\cite{bollen2011} reported that aggregate
mood extracted from social media carries predictive content for index movements. Da, Engelberg
and Gao~\cite{da2011attention} demonstrated that retail attention itself, measured through search
behaviour, moves prices, and Barber and Odean~\cite{barber2008} documented the mechanism by which
attention-grabbing news drives individual buying. This body of work establishes the premise on
which the sentiment channel of this thesis rests; what it does not establish is how much of the
information survives into an out-of-sample forecast, which is the question
Chapter~\ref{ch:llm} measures.

The evolution from lexicon methods through supervised learning to transformer language models
was described technically in Section~2.3. Its research significance is that each generation
resolved the failure of the previous one. Loughran and McDonald~\cite{loughran2011} established
that general-purpose lexicons misclassify financial language badly enough to require
domain-specific alternatives, which is the finding that made financial sentiment analysis a
distinct field rather than an application of general sentiment analysis;
VADER~\cite{hutto2014} remains the reference rule-based implementation for social text.
Supervised machine learning improved accuracy at the cost of annotation and generalised poorly
across market conditions~\cite{rouf2021}. Transformer architectures~\cite{vaswani2017} displaced
both, and domain adaptation followed: FinBERT~\cite{araci2019} adapted bidirectional encoders to
financial language; BloombergGPT~\cite{wu2023bloomberggpt} and the FinGPT
family~\cite{yang2023fingpt,wang2024fingptbench} extended the approach to generative reasoning
over financial documents. Surveys~\cite{li2024llmsurvey,todd2024,bahoo2024} converge on the
finding that contextual language models extract market-relevant meaning that lexicon methods
miss.

\subsection{Applications and their common evaluation weakness}

Applications have followed quickly and cover a wide range of markets and instruments. Fatouros
et al.~\cite{fatouros2024} showed that conversational language models outperform established
financial sentiment baselines. Wilksch and Abramova~\cite{wilksch2024} built a dedicated model
for sentiment in financial tweets. Haryono et al.~\cite{haryono2023} coupled a transformer news
encoder with a gated recurrent unit for technical data and later extended the approach to
generative aspect-based sentiment~\cite{haryono2024}. Ho and Huang~\cite{ho2021} combined
sentiment with candlestick chart representations. Jung and Jang~\cite{jung2024} preserved
quantitative context by masking numerically significant tokens. Ónózó et
al.~\cite{onozo2024} applied language models to news analysis and macroeconomic nowcasting.
Vargas et al.~\cite{vargas2022} fused sentiment with recurrent price models on Brazilian
equities, Wang and Chen~\cite{wang2024spcm} built a sentiment-based recommendation system,
Shang~\cite{shang2025} combined high-dimensional indicators with social media dynamics, Elahi
and Taghvaei~\cite{elahi2024} combined financial data with news articles using language models
directly, Officioso et al.~\cite{officioso2026} coupled financial news sentiment with market
data in a language-model-based predictor, and Mishra et al.~\cite{mishra2025} and Liu et
al.~\cite{liu2024llmreview} surveyed the resulting body of work.

What this literature establishes is that language models improve sentiment
\emph{measurement}. What it establishes far less clearly is how much of that improvement
survives into \emph{forecast quality}. Three evaluation weaknesses recur. First, the reported
gain is frequently stated without the baseline against which it should be judged: an accuracy
of 55 per cent is impressive against a balanced base rate and unremarkable against a class
imbalance of the same magnitude. Second, temporal alignment is often described loosely, and the
attribution of a post-close publication to the session that has already ended is a subtle form
of look-ahead that inflates results without leaving a trace. Third, evaluation stops at the
classification metric; the question of whether a measured increment in accuracy is large enough
to survive transaction costs is rarely posed.

\subsection{Market-derived sentiment as a distinct construct}

A second and older family of sentiment measures reads sentiment off the market rather than off
text. Baker and Wurgler~\cite{bakerwurgler2007} construct investor sentiment from trading
activity and dispersion, on the argument that participation reveals what surveys merely report.
Singsiri and Yokrattanasak~\cite{singsiri2026} build allocation signals from money-flow
indicators on exactly this reasoning, and Naidoo et al.~\cite{naidoo2025} model investor
sentiment and market volatility as a joint object using a mixed-frequency conditional
heteroskedasticity framework.

This family matters to the present thesis for a specific structural reason. Instrument-level,
time-stamped news corpora of the depth required for a cross-sectional study do not exist for
broad index instruments such as exchange-traded funds. A cross-sectional allocation study over
such instruments therefore cannot use textual sentiment and must use a market-derived
construct, which is what Chapter~\ref{ch:allocation} does. The consequence --- that a null
result on the market-derived composite is a null about \emph{that proxy}, not about textual
sentiment in general --- is recorded there and is one of the more important qualifications in
the thesis.

\section{Explainability, Causality and Uncertainty}

\subsection{Post-hoc attribution and its limit}

Opacity is a recognised obstacle to the adoption of learned models in finance. Koa et
al.~\cite{koa2024} used self-reflective large language models to generate natural-language
explanations alongside stock predictions, and the general-purpose attribution methods of
Lundberg and Lee~\cite{lundberg2017} and Ribeiro et al.~\cite{ribeiro2016} are now standard
instruments for post-hoc interpretation of fitted models.

The limitation is structural rather than a matter of method quality. Post-hoc attribution
explains a model: it reports which inputs the fitted function relied upon in reaching a
particular output. It does not establish that the relationships the model has learned are
anything more than correlations that happened to hold in the training window. A feature that
co-moves with the target for spurious reasons receives an attribution exactly as a genuine
driver would, and an audience shown that attribution has been given a faithful account of the
model and a potentially misleading account of the market.

This concern is not idiosyncratic to finance. The comprehensive survey of explainable artificial
intelligence by Barredo Arrieta et al.~\cite{arrieta2020} distinguishes transparency by design
from post-hoc interpretability and notes that the two answer different questions; Doshi-Velez and
Kim~\cite{doshivelez2017} argue that an interpretability claim is empty without an explicit
statement of what the explanation is for and who must act on it; and Rudin~\cite{rudin2019} makes
the stronger case that for high-stakes decisions an intrinsically interpretable model should be
preferred to an explained black box. The position taken in this thesis is a version of Rudin's
applied to the feature set rather than to the model class: constrain \emph{what may enter} the
model by a structural criterion, and the resulting system becomes auditable regardless of how the
fitting is subsequently done.

\subsection{Causal formulations}

Causal formulations address this directly. Xu, Cheng and Lee~\cite{xu2026causal} showed that
causal feature selection is more robust out of sample than correlation-based feature
engineering, separating genuine causal drivers from spurious associations. Their causal graph
is static, however, and cannot represent the structural breaks produced by policy
announcements --- a limitation that is acute in the Indian setting, where budget and monetary
policy events reset inter-instrument relationships within hours.

Liu et al.~\cite{liu2026dynamic} introduced a dynamic-causal lens that encodes stock
interactions as symbolic movement patterns and adapts to changing conditions. The formulation
is the closest in the literature to what the Indian setting requires, and its cost is
computational: the approach scales quadratically in the number of instruments, which is
prohibitive for a fifty-constituent index. Yu and Liu~\cite{yu2025causalseqgnn} coupled factor
investing with causal graph neural networks through a fuzzy cognitive map, demonstrating that
multi-factor causal structure improves predictive power; their factors are Fama--French factors
calibrated to United States equities, which do not reflect Indian idiosyncrasies such as the
weight of public-sector banks or the seasonal influence of the monsoon on agricultural
commodities and the stocks exposed to them. Luo et al.~\cite{luo2025magicnet} added memory to a
causal interaction network, with the attendant risk of over-fitting to historical crises whose
character does not repeat; and Xing et al.~\cite{xing2024vague} modelled the vagueness of
inter-stock relations, recognising that inter-stock relationships lie on a continuum rather than
being crisp, without supplying calibration guidance for an emerging-market setting.

\subsection{Multi-granularity temporal structure}

Alongside the causal line, multi-granularity architectures demonstrated that financial series
carry structure at several time scales at once. Zhu et al.~\cite{zhu2025multigran} processed
minute-, day- and week-level data jointly and showed that cross-granularity information
transfer improves trend prediction; their graph construction assumes homogeneity across
instruments, which fails in a universe where highly liquid constituents behave differently from
thinly traded ones. Wang et al.~\cite{wang2025mtmd} added a debiasing layer to multi-scale
temporal memory learning that reduces historical bias in trend forecasting.

The debiasing layer carries an assumption that deserves emphasis, because
Chapter~\ref{ch:causal} reports it failing empirically. The correction assumes \emph{stationary}
noise properties, so that a bias estimated on past data remains valid. Indian market volatility
is not stationary: it clusters, responds asymmetrically to negative returns, and changes
character discretely at policy events. A debiasing layer built on a stationarity assumption is
therefore least reliable at exactly the moments when reliability matters most.

\subsection{Relational learning across the cross-section}

Relational formulations extend the same idea to the cross-section. Patel, Jariwala and
Chattopadhyay~\cite{patel2024hybrid} treat securities as nodes in a learned relational graph and
show that relational modelling improves directional accuracy over independent treatment; their
model predicts without dynamic risk adjustment, which is a serious limitation for a market with
intraday reversals. Liu et al.~\cite{liu2024heterogeneous} model mutual interplay through
dual-dynamic attention across cross-sectional and temporal dimensions, with scalability to large
universes not fully addressed. Gao et al.~\cite{gao2025relational} select stocks through
probabilistic state-space learning, assuming linear Gaussian transitions that are strained by
the nonlinear dynamics of election periods and policy announcements.

\subsection{Uncertainty as a first-class quantity}

A related and less examined issue is that a signal's reliability is itself uncertain. Most
fusion architectures pass point predictions forward without any statement of confidence, so a
source that has recently become noisy continues to be weighted as though it had not. In Bayesian
machine learning a prediction is incomplete without an estimate of its uncertainty, and the
distinction between \emph{aleatoric} uncertainty --- irreducible noise in the environment ---
and \emph{epistemic} uncertainty --- the model's ignorance about regions it has not observed ---
is well established~\cite{gal2016dropout,kendall2017}. Khuwaja, Khowaja and
Dev~\cite{khuwaja2023} address heterogeneity in financial data through adversarial learning
without making per-source confidence explicit. Across the fusion literature surveyed in
Section~\ref{sec:fusionsurvey}, explicit, time-varying uncertainty quantification per source is
absent.

\section{From Prediction to Decision: Allocation and Validation}

\subsection{Portfolio construction and the estimation problem}

Portfolio construction has its own long literature. Markowitz~\cite{markowitz1952} established
the mean-variance formulation, whose extreme sensitivity to estimated inputs motivated both
shrinkage estimators~\cite{ledoit2004} and the demonstration by DeMiguel, Garlappi and
Uppal~\cite{demiguel2009} that naive equal weighting can outperform optimisation that ignores
estimation error. That demonstration is the reason equal weighting is the primary benchmark in
Chapter~\ref{ch:allocation}: a framework that cannot beat $1/N$ has not earned its complexity.
Risk-based construction, of which the hierarchical approach of López de
Prado~\cite{lopezdeprado2016} is representative, avoids the estimation of expected returns
entirely by forgoing signal information --- a defensible trade that nonetheless defines away the
question this thesis asks.

\subsection{Regime-aware and risk-managed allocation}

Regime-switching models following Hamilton~\cite{hamilton1989} showed that financial series are
not homogeneous through time. Ang and Bekaert~\cite{ang2002} demonstrated that accounting for
regime shifts materially changes optimal allocations, and Kritzman, Page and
Turkington~\cite{kritzman2012} made the practical case for regime-aware dynamic strategies. Luo
and Mulvey~\cite{luo2026regime} extend the line with dual-regime signals and regime-dependent
asset selection. Volatility-managed portfolios~\cite{moreira2017} scale total risk by recent
volatility and improve risk-adjusted returns, but say nothing about composition: the mechanism
is orthogonal to which assets are held.

Kumar and Yadav~\cite{kumar2024rdmdea} proposed credibilistic portfolio selection under a range
directional measure formulation, handling positive and negative returns without assuming
normality, which is a genuine advance for return distributions that are neither symmetric nor
thin-tailed. The formulation is not coupled to a dynamic signal, so rebalancing occurs at fixed
intervals rather than in response to a change in the underlying relationships. Gu, Du and
Wang~\cite{gu2025ragic} generate risk-aware prediction intervals rather than point forecasts,
requiring calibration that is not supplied for emerging-market instruments; Zhu et
al.~\cite{zhu2025spin} build sparse portfolios driven by irregular news streams, using
English-language sources only, which for an Indian deployment omits a substantial share of the
information actually reaching retail participants.

\subsection{Evaluation practice and the credibility of a reported gain}

Evaluation practice is the final concern and, in the judgement of this thesis, the most
consequential. Risk-adjusted performance is still summarised by the ratio introduced by
Sharpe~\cite{sharpe1966}, and Gu, Kelly and Xiu~\cite{gu2020} provide the benchmark
machine-learning study in asset pricing, noting that measured predictability is small. Harvey
and Liu~\cite{harvey2015} showed that conventional significance thresholds are inadequate once
the number of strategies examined is taken into account, and Bailey and López de
Prado~\cite{bailey2014dsr} formalised the deflated Sharpe ratio, which adjusts for the number of
trials, the non-normality of returns and the sample length. Valls and Steurer~\cite{valls2026}
document in detail how volatility-based strategies overfit and what the resulting reported
performance is worth. Robust inference under serial dependence requires
corrections~\cite{newey1987}, resampling schemes~\cite{politis1994} and forecast-comparison
tests~\cite{diebold1995}.

The practical consequence adopted throughout this thesis is a disclosure discipline rather than
a test: a reported gain is credible only when the specification search behind it is reported
alongside it. Reporting the best cell of a grid and omitting the grid is not a presentational
choice but a misstatement of the evidence.

\section{Consolidated Review and Positioning}

Table~\ref{tab:lit} consolidates the representative work reviewed above, stating for each the
contribution that is relevant to this thesis and the limitation that this thesis must work
around. The final column is not a criticism of the cited work, which is in each case addressing
the question it set itself; it records the reason the work cannot be adopted unmodified here.

\begin{table}[!htb]
\centering
\caption[Consolidated review of representative work]{Consolidated review of representative work,
with the limitation relevant to this research}
\label{tab:lit}
\tabfonts
\begin{tabular}{@{}L{2.5cm}L{3.3cm}L{3.9cm}L{4.1cm}@{}}
\toprule
\textbf{Reference} & \textbf{Technique} & \textbf{Contribution relevant here} &
\textbf{Limitation to be addressed} \\
\midrule
Snášel et al.~\cite{snasel2024} & Generalised multi-source fusion for stock selection & Formal
basis for combining numerical and textual streams & Fusion weights are not conditioned on the
prevailing market state \\
Long et al.~\cite{long2024}; Nti et al.~\cite{nti2021} & Multi-view and multi-source deep
fusion & Views carry non-overlapping information; frequency alignment is essential & Developed
and tested outside emerging-market conditions \\
Fu and Zhang~\cite{fu2024} & Multi-source market sentiment with price data & News and social
sentiment behave differently and should be separate streams & Combination weights are fixed
across regimes \\
Haryono et al.~\cite{haryono2023}; Mu et al.~\cite{mu2023investor} & Sentiment fused with
technical indicators & Gated and optimised fusion of text and price streams & Gate learned
globally; sentiment quality assumed rather than validated \\
Farimani et al.~\cite{farimani2024} & Adaptive multimodal weighting & Moves toward time-varying
source importance & No explicit probabilistic confidence per source \\
Araci~\cite{araci2019}; Wu et al.~\cite{wu2023bloomberggpt} & Domain-adapted language models &
Contextual sentiment extraction from financial text & Gain in sentiment quality is not traced
through to decision quality \\
Jung and Jang~\cite{jung2024} & Numerical-masking adaptation & Preserves quantitative context
that generic pre-processing destroys & Evaluated on measurement quality, not on forecast value
\\
Koa et al.~\cite{koa2024}; Lundberg and Lee~\cite{lundberg2017} & LLM explanations; post-hoc
attribution & Natural-language and feature-level justification & Explains the model, not the
causal structure \\
Xu et al.~\cite{xu2026causal} & Causal perspective on stock prediction & Causal features are
more robust out of sample & Static graph cannot absorb policy shocks \\
Liu et al.~\cite{liu2026dynamic} & Dynamic-causal graph from symbolic movement patterns &
Structure adapts to changing conditions & Scales quadratically; no sector-specific adaptation \\
Zhu et al.~\cite{zhu2025multigran} & Multi-granularity graph augmented learning & Several time
scales processed jointly & Assumes homogeneity across instruments \\
Wang et al.~\cite{wang2025mtmd} & Multi-scale temporal memory with debiasing & Multiple horizons
processed jointly & Debiasing assumes stationary noise \\
Kumar and Yadav~\cite{kumar2024rdmdea} & RDM-DEA credibilistic portfolio selection &
Optimisation without a normality assumption & Not coupled to a dynamic signal for rebalancing \\
Moreira and Muir~\cite{moreira2017} & Volatility-managed portfolios & Total risk scaled by
recent volatility & Does not address portfolio composition \\
Zhu et al.~\cite{zhu2025spin} & Sparse portfolios from irregular news & Sparse allocation
responsive to news flow & English-only sources; omits regional-language sentiment \\
DeMiguel et al.~\cite{demiguel2009}; Harvey and Liu~\cite{harvey2015} & Naive diversification;
backtesting under multiple testing & Estimation error can defeat optimisation; gains must
account for specification search & Sets the benchmark any framework must beat and requires every
specification to be disclosed \\
\bottomrule
\end{tabular}
\end{table}

\section{Research Gaps}

The review supports five gaps. Figure~\ref{fig:taxonomy} positions them against the state of the
art along the three axes of fusion strategy, signal extraction and decision layer.

\begin{enumerate}[label=\textbf{G\arabic*.},leftmargin=2.8em]
\item \textbf{Static fusion of dynamically reliable sources.} Existing frameworks fuse
financial, textual and macroeconomic streams using fixed weights, globally learned attention or
simple averaging. As Table~\ref{tab:fusiontax} records, none of the surveyed studies conditions
the contribution of a source on a measured, time-varying estimate of that source's own
predictive reliability. A stream that has become noisy therefore continues to influence the
forecast at its historical weight until the next retraining cycle.

\item \textbf{Incomplete coverage of the Indian information environment.} Most multi-source
frameworks are designed and validated on developed markets. Frameworks that jointly incorporate
financial indicators, news sentiment, social media signals and Indian macroeconomic variables
--- policy rate, inflation, industrial production, currency --- under a single architecture and
a strict temporal protocol are rare, and the calibrations imported from developed markets
(Fama--French factors, crisis memory, homogeneity assumptions) do not transfer cleanly.

\item \textbf{Unverified translation from language-model sentiment to decision quality.}
Language models demonstrably improve sentiment measurement. Whether that improvement survives
temporal alignment, fusion with technical indicators and honest out-of-sample evaluation at a
single-name horizon has not been established, and reported gains are seldom accompanied by the
baseline against which they should be judged or by an assessment of whether the increment is
large enough to act on.

\item \textbf{Correlational rather than mechanistic transparency.} Post-hoc attribution explains
a fitted model. It does not distinguish a genuine driver from a spurious association, and it
does not describe how a sentiment shock propagates into a price movement over time. A
transparency mechanism grounded in directed causal structure and in multi-scale temporal
dynamics is required, together with a direct measurement of how much noise the sentiment channel
actually carries --- as distinct from an assumption that it carries little.

\item \textbf{Prediction reported in place of decision.} Statistical accuracy remains the
dominant success criterion, and economic value is reported inconsistently. Where allocation is
addressed it is often heuristic, and the contribution of individual architectural layers is
almost never isolated, so architectures accumulate components whose value has never been tested.
Cross-market transfer is asserted more often than it is demonstrated, and the conditions under
which a framework should be expected to transfer are rarely stated in advance.
\end{enumerate}

\begin{figure}[!htb]
\centering
\scalebox{0.86}{\begin{tikzpicture}[node distance=4mm]
\node[blkd,minimum width=34mm] (root) {Multi-source financial\\prediction and decision};

\node[blk,below left=6mm and 24mm of root,minimum width=27mm,minimum height=6mm] (a) {Fusion strategy};
\node[blk,below=6mm of root,minimum width=27mm,xshift=-2mm,minimum height=6mm] (b) {Signal extraction};
\node[blk,below right=6mm and 22mm of root,minimum width=27mm,minimum height=6mm] (c) {Decision layer};

\node[blk,below=5mm of a,minimum width=29mm,minimum height=6mm] (a1)
  {early / late / attention\\\textit{fixed or global weights}};
\node[blk,below=5mm of b,minimum width=29mm,minimum height=6mm] (b1)
  {lexicon $\to$ ML $\to$ LLM\\\textit{measurement improves}};
\node[blk,below=5mm of c,minimum width=29mm,minimum height=6mm] (c1)
  {accuracy reported\\\textit{allocation heuristic}};

\node[blka,below=5mm of a1,minimum width=29mm,minimum height=9mm] (a2)
  {\textbf{G1, G2}\\uncertainty-weighted,\\market-specific fusion};
\node[blka,below=5mm of b1,minimum width=29mm,minimum height=9mm] (b2)
  {\textbf{G3, G4}\\validated sentiment,\\causal transparency};
\node[blka,below=5mm of c1,minimum width=29mm,minimum height=9mm] (c2)
  {\textbf{G5}\\ablated allocation,\\cross-market test};

\draw[ar] (root) -- (a); \draw[ar] (root) -- (b); \draw[ar] (root) -- (c);
\draw[ar] (a) -- (a1); \draw[ar] (b) -- (b1); \draw[ar] (c) -- (c1);
\draw[arb] (a1) -- (a2); \draw[arb] (b1) -- (b2); \draw[arb] (c1) -- (c2);
\node[lbl,rotate=90,anchor=center] at ([xshift=-9mm]a1.west) {state of the art};
\node[lbl,rotate=90,anchor=center] at ([xshift=-9mm]a2.west) {this research};
\end{tikzpicture}}
\caption[Taxonomy of the reviewed field and the five research gaps]{Taxonomy of the reviewed
field and the location of the five research gaps addressed by this work.}
\label{fig:taxonomy}
\end{figure}

\section{Problem Statement and Research Objectives}

\subsection{Problem statement}

To design, develop and validate an integrated decision-support framework that fuses
heterogeneous financial information --- market metrics, news and social media sentiment
extracted using large language models, and macroeconomic indicators --- into interpretable and
reliability-aware predictions, and that converts those predictions into executable,
risk-controlled stock market decisions whose performance is demonstrated across market states
and across markets. The framework must address the limitations of existing approaches, namely
static source weighting under non-stationary reliability, opaque and correlational reasoning,
unquantified sentiment noise, and the routine substitution of statistical accuracy for economic
value.

\subsection{Research objectives}

\begin{enumerate}[label=\textbf{RO\arabic*.},leftmargin=2.8em]
\item To investigate and develop a comprehensive model that integrates diverse data sources ---
such as financial metrics, news articles, social media sentiment and macroeconomic factors ---
aiming to enhance the accuracy of stock market predictions by providing a holistic analysis.

\item To design a hybrid framework that combines large language models, sentiment analysis
techniques and conventional financial indicators, focusing on optimising deep learning
architectures for real-time stock market predictions.

\item To establish and evaluate an explainable artificial intelligence framework capable of
delivering transparent and interpretable predictions in stock selection, analysing the impact of
sentiment data noise and exploring the temporal dynamics between sentiment shifts and stock
price movements.

\item To evaluate and validate the result and practical application of the proposed model across
different stock markets, exploring transfer learning techniques for cross-market applications
and validating the model against benchmark datasets and real-world market data.
\end{enumerate}

\subsection{Alignment of gaps, objectives and contributions}

Figure~\ref{fig:gapmap} traces each gap to the objective that addresses it and to the
contribution reported in the corresponding chapter. The mapping is many-to-many by design: G1
and G2 are jointly answered by the fusion work of RO1 across Chapters~\ref{ch:fusion}
and~\ref{ch:dynfusion}, while G5 requires evidence from both the decision layer of RO4 and the
transparency instruments of RO3.

\begin{figure}[!htb]
\centering
\scalebox{0.86}{\begin{tikzpicture}[node distance=3mm]
\node[blkd,minimum width=32mm,minimum height=6mm] (g1) {\textbf{G1} static fusion weights};
\node[blkd,below=3mm of g1,minimum width=32mm,minimum height=6mm] (g2) {\textbf{G2} Indian information environment};
\node[blkd,below=3mm of g2,minimum width=32mm,minimum height=6mm] (g3) {\textbf{G3} sentiment $\to$ decision unverified};
\node[blkd,below=3mm of g3,minimum width=32mm,minimum height=6mm] (g4) {\textbf{G4} correlational transparency};
\node[blkd,below=3mm of g4,minimum width=32mm,minimum height=6mm] (g5) {\textbf{G5} prediction reported as decision};

\node[blkg,right=26mm of g1,minimum width=28mm,minimum height=6mm] (o1) {\textbf{RO1} integrated multi-source model};
\node[blkg,below=5mm of o1,minimum width=28mm,minimum height=6mm] (o2) {\textbf{RO2} hybrid LLM + indicator framework};
\node[blkg,below=5mm of o2,minimum width=28mm,minimum height=6mm] (o3) {\textbf{RO3} explainable and causal framework};
\node[blkg,below=5mm of o3,minimum width=28mm,minimum height=6mm] (o4) {\textbf{RO4} validation and cross-market transfer};

\node[blka,right=22mm of o1,minimum width=30mm,minimum height=6mm] (c1) {\textbf{Ch.\ 4--5} MSFN and Bayesian fusion};
\node[blka,below=5mm of c1,minimum width=30mm,minimum height=6mm] (c2) {\textbf{Ch.\ 6} LLM--indicator pipeline};
\node[blka,below=5mm of c2,minimum width=30mm,minimum height=6mm] (c3) {\textbf{Ch.\ 7} DCHF and sentiment diagnostics};
\node[blka,below=5mm of c3,minimum width=30mm,minimum height=6mm] (c4) {\textbf{Ch.\ 8} decision-theoretic allocation};

\draw[ar] (g1) -- (o1); \draw[ar] (g2) -- (o1);
\draw[ar] (g3) -- (o2); \draw[ar] (g4) -- (o3);
\draw[ar] (g5) -- (o4); \draw[arb] (g5) -- (o3);
\draw[ar] (o1) -- (c1); \draw[ar] (o2) -- (c2);
\draw[ar] (o3) -- (c3); \draw[ar] (o4) -- (c4);
\node[lbl,above=1.5mm of g1] {Research gaps};
\node[lbl,above=1.5mm of o1] {Research objectives};
\node[lbl,above=1.5mm of c1] {Chapters and contributions};
\end{tikzpicture}}
\caption[Alignment of research gaps, objectives and thesis chapters]{Alignment of research
gaps, research objectives and the chapters in which the corresponding contributions are
developed. Dashed links denote partial rather than primary coverage.}
\label{fig:gapmap}
\end{figure}

\section{Summary}

The literature establishes that multi-source fusion improves financial prediction, that language
models improve sentiment measurement, that causal feature selection is more robust out of sample
than correlational selection, and that regime awareness changes optimal allocations. It leaves
four things unresolved that this thesis addresses: fusion weights are learned once rather than
measured continuously; the Indian information environment is under-represented and the
calibrations imported from developed markets do not transfer; the translation from improved
sentiment measurement to improved decisions is asserted rather than demonstrated; and the
decision layer is evaluated as an aggregate, so that the contribution of individual architectural
components is unknown. The next five chapters address these in turn.
